\documentclass[a4paper,12pt]{article}
\usepackage[utf8]{inputenc}
\usepackage[T1]{fontenc}
\usepackage[italian]{babel}
\usepackage{listings}
\usepackage{xcolor}
\usepackage{hyperref}
\usepackage[margin=3cm]{geometry}

\title{Documentazione Progetto Sistemi Operativi: PandOSsh Fase 2}
\author{Luca Argentino, Milo Disalvatore, Matteo Mazzetti, Zeno Maccari}
\begin{document}

\maketitle

\section{Inizializzazione (initial.c)}
...

\vspace{0.5cm}
\subsection{Strutture Dati Globali}
Le variabili globali definite nel modulo sono:
\begin{itemize}
    \item ...
\end{itemize}

\vspace{0.5cm}
\subsection{Descrizione delle Funzioni}

\subsubsection{\texttt{pcb\_t* allocFirstPcb()}}
...

\subsubsection{\texttt{void initVector()}}
...

\subsubsection{\texttt{int main()}}
...

\section{Scheduler (scheduler.c)}
...

\vspace{0.5cm}
\subsection{Descrizione delle Funzioni}

\subsubsection{\texttt{void dispatch()}}
...

\section{Gestione Interrupt (interrupts.c)}
Il modulo \texttt{interrupts.c} gestisce le eccezioni di interrupt sia dei timer (Processor Local Timer e Interval Timer) sia dei dispositivi hardware. I dispositivi gestiti sono Disk devices, Flash devices, Network devices, Printer devices e Terminal devices(che possono essere di tipo transmitter o reciever).

\vspace{0.5cm}
\subsection{Descrizione delle Funzioni}

\subsubsection{\texttt{void pltInterruptHandler()}}
\begin{enumerate}
    \item Aumenta il PLT (timer locale del processore) di un quanto di 5ms.
    \item Aggiorna il tempo di utilizzo del processore del processo corrente \texttt{currProc}.
    \item Memorizza lo stato attuale del processore all'istante esatto in cui è avvenuto un interrupt e lo memorizza nel campo \texttt{p\_s} del PCB del processo attuale.
    \item Inserisce il processo interrotto nella \texttt{readyQueue}.
    \item Invoca la funzione \texttt{dispatch()} dello scheduler. 
\end{enumerate}

\subsubsection{\texttt{void itInterruptHandler()}}
\begin{enumerate}
    \item Ricarica l'Interval Timer con un nuovo time slice pari a 100ms.
    \item Rilascia tutti i processi in attesa di un tick di timer e li inserisce nella \texttt{readyQueue}, diminuendo in numero di processi soft-bloccati (in attesa di eseguire un operazione I/O).
    \item Resetta il valore dell'ultimo semaforo.
    \item Salva lo stato attuale del processore nel momento in cui è avvenuto un interrupt e lo memorizza nel campo \texttt{p\_s} del PCB del processo attuale.
    \item Richiama l'unica procedura dello scheduler: \texttt{dispatch()}. 
\end{enumerate}

\subsubsection{\texttt{void deviceInterruptHandler(int excCode)}}
\begin{enumerate}
    \item Mappa l'\texttt{excCode} che ottiene dal parametro a una delle righe corrispondenti agli interrupt hardware che possono essere: \begin{itemize}
        \item Line 3: Disk
        \item Line 4: Flash
        \item Line 5: Ethernet
        \item Line 6: Printer
        \item Line 7: Terminal
    \end{itemize} 
    \item Calcola l'indirizzo di memoria del device che genera l'interrupt ed esegue operazioni bit a bit per determinare quale dei vari device elencati sopra ha effettivamente generato l'interrupt.
    \item Calcola l'indirizzo fisico di memoria dei registri di tale device in \texttt{devAddr} e ottiene l'indice del semaforo corrispondente.
    \item Controlla se l'interrupt è generato da un terminale e in tal caso verifica se esso è di tipo transmitter o reciever, leggendo di conseguenza l'appropriato registro di stato.
    \item Esegue un operazione di tipo Verhoog sul semaforo di quel device: sbloccando il primo processo in attesa o facendo sì, se la coda è vuota, che il prossimo non si blocchi. Se un prcesso era in attesa di questo I/O: \begin{itemize}
        \item Memorizziamo il codice di stato di quel processo nel registro \texttt{a0} che corrisponde a \texttt{p\_s->gpr[24]}.
        \item Decrementa il \texttt{softBlock\_count}, ossia il contatore dei processi che aspettano un'I/O.
        \item Restituisce il controllo al processo interrotto
    \end{itemize}
    Altrimenti viene richiamato lo scheduler.
\end{enumerate}

\subsubsection{\texttt{void interruptHandler()}}
\begin{enumerate}
    \item Esegue un \texttt{AND} bitwise tra l'\texttt{excState} e un'apposita maschera e determina la tipologia dell'interrupt.
    \item  Se era associato al Processor Local Timer richiama \texttt{pltInterruptHandler()}.
    \item Se era associato all'Interval Timer invoca \texttt{itInterruptHandler()}.
    \item Se invece era associato ad un device passa il controllo a \texttt{deviceInterruptHandler()}.
\end{enumerate}

\section{Eccezioni (exceptions.c)}
...

\vspace{0.5cm}
\subsection{Descrizione delle Funzioni}

\subsubsection{\texttt{void exceptionHandler()}}
...

\subsubsection{\texttt{void passUpOrDie(int)}}
...

\subsubsection{\texttt{void exc\_tlbHandler()}}
...

\subsubsection{\texttt{void exc\_trapHandler()}}
...

\subsubsection{\texttt{void syscallHandler()}}
...

\vspace{0.5cm}
\section{Chiamate di Sistema (syscall.c)}
...

\vspace{0.5cm}
\subsection{Descrizione delle Funzioni}

\subsubsection{\texttt{int nsys1(int a1, int a2, int a3)}}
...

\subsubsection{\texttt{void nsys2(int a1, state\_t *excState)}}
...

\subsubsection{\texttt{void nsys3(int a1, state\_t *excState)}}
...

\subsubsection{\texttt{void nsys4(int a1, state\_t *excState)}}
...

\subsubsection{\texttt{void nsys5(int a1, int a2, state\_t *excState)}}
...

\subsubsection{\texttt{void nsys6(state\_t *excState)}}
...

\subsubsection{\texttt{void nsys7(state\_t *excState)}}
...

\subsubsection{\texttt{void nsys8(state\_t *excState)}}
...

\subsubsection{\texttt{void nsys9(int a1, state\_t *excState)}}
...

\subsubsection{\texttt{void nsys10(state\_t *excState)}}
...

\end{document}
